\documentclass[12pt, twoside]{article}
\input{../preamble}

\fancyhead[LO]{La Scuola d'Italia / Dr.\ Huson / IB Math 11 \\* 16 September 2026}
\fancyhead[RO]{\fbox{\begin{minipage}{6cm}\footnotesize Nickname:\\[0.5cm]\end{minipage}}}
\fancyhead[LE]{\fbox{\begin{minipage}{6cm}\footnotesize Nickname:\\[0.5cm]\end{minipage}}}
\fancyfoot[C]{\thepage}

\begin{document}
\subsubsection*{1.4 Pre-Quiz: Arithmetic Sequences}

\emph{Homework. Due Monday 21 September, at the start of class. Write your
answers in the boxes; a calculator is allowed. Several parts ask you to write
a sentence --- those sentences are the point of the assignment.}

\begin{enumerate}[itemsep=0.15cm]

%--- Problem 1: Arithmetic vs non-arithmetic; the constant-difference test [9 marks] ---
%--- Source: Original (freshly authored, AA SL 1.2) ---
\item[1.]

Here are five sequences.

\begin{center}
\begin{tabular}{rl}
(i) & $6,\; 10,\; 14,\; 18,\; 22,\; \ldots$ \\[2pt]
(ii) & $3,\; 6,\; 12,\; 24,\; 48,\; \ldots$ \\[2pt]
(iii) & $1,\; 4,\; 9,\; 16,\; 25,\; \ldots$ \\[2pt]
(iv) & $20,\; 14,\; 8,\; 2,\; -4,\; \ldots$ \\[2pt]
(v) & $5,\; 5,\; 5,\; 5,\; 5,\; \ldots$ \\
\end{tabular}
\end{center}

\begin{enumerate}[itemsep=0.1cm]
  \item For each sequence, write down the three differences
  $u_{2} - u_{1}$, $u_{3} - u_{2}$ and $u_{4} - u_{3}$.
  \item State which of the five sequences are arithmetic, and write down
  the common difference $d$ of each one.
  \item For each sequence that is \emph{not} arithmetic, describe its rule
  in one sentence.
  \item Complete the sentence: ``A sequence is arithmetic when \ldots''
\end{enumerate}

\begin{tikzpicture}
  \draw (-0.6,0) rectangle (14.6,11.0);
  \foreach \y in {1,2,...,10} \draw[dotted] (0,\y)--(14.1,\y);
\end{tikzpicture}

\newpage

%--- Problem 2: Sequence vs series; u_n vs S_n notation [9 marks] ---
%--- Source: Original (freshly authored, AA SL 1.2) ---
\item[2.]

Two words in this unit sound alike but mean different things:
\emph{sequence} and \emph{series}.

\begin{enumerate}[itemsep=0.1cm]
  \item Explain the difference between a sequence and a series, in your own
  words.
  \item Consider the sequence $4,\; 9,\; 14,\; 19,\; 24$. Write down the
  series formed from these five terms, and find its value.
  \item For that sequence, state what $u_{5}$ means and what $S_{5}$ means,
  and write down the value of each.
  \item State whether each of the following is a sequence or a series:

  \hspace*{1em} $3,\; 7,\; 11,\; 15$ \hfill $3 + 7 + 11 + 15$ \hfill
  $u_{n} = 2n + 1$ \hfill $S_{20} = 430$ \hspace*{1em}
  \item Explain why $S_{5} - S_{4} = u_{5}$.
\end{enumerate}

\begin{tikzpicture}
  \draw (-0.6,0) rectangle (14.6,14.5);
  \foreach \y in {1,2,...,13} \draw[dotted] (0,\y)--(14.1,\y);
\end{tikzpicture}

\newpage

%--- Problem 3: The Gauss story; pairing and the S_n = n/2 (u_1 + u_n) formula [9 marks] ---
%--- Source: Original (freshly authored, AA SL 1.2) ---
\item[3.]

When Carl Friedrich Gauss was a schoolboy, his teacher asked the class to
add all the whole numbers from $1$ to $100$. Gauss wrote down the answer
almost at once.

\begin{enumerate}[itemsep=0.1cm]
  \item Find the value of $1 + 100$, of $2 + 99$, and of $3 + 98$. Write
  down what you notice.
  \item How many pairs like these are there in $1 + 2 + 3 + \cdots + 100$?
  Hence find the total.
  \item The sum $1 + 2 + 3 + \cdots + 100$ is an arithmetic series. Write
  down $u_{1}$, $d$ and $u_{100}$, then use
  $S_{n} = \dfrac{n}{2}\bigl(u_{1} + u_{n}\bigr)$ to check your answer
  to (b).
  \item Explain, in one or two sentences, why Gauss's pairing and the
  formula are the same idea. Where does the $\dfrac{n}{2}$ come from?
  \item Use Gauss's method to find $1 + 2 + 3 + \cdots + 500$.
\end{enumerate}

\begin{tikzpicture}
  \draw (-0.6,0) rectangle (14.6,14.0);
  \foreach \y in {1,2,...,13} \draw[dotted] (0,\y)--(14.1,\y);
\end{tikzpicture}

\newpage

%--- Problem 4: Fibonacci; recursive vs explicit rules [9 marks] ---
%--- Source: Original (freshly authored, AA SL 1.2) ---
\item[4.]

The \emph{Fibonacci} sequence begins
$$1,\; 1,\; 2,\; 3,\; 5,\; 8,\; 13,\; \ldots$$
Each term after the second is the sum of the two terms before it.

\begin{enumerate}[itemsep=0.1cm]
  \item Write down the next three terms.
  \item Show that the Fibonacci sequence is not arithmetic.
  \item The rule $u_{n+2} = u_{n+1} + u_{n}$, with $u_{1} = u_{2} = 1$, is
  called a \textbf{recursive} rule. Explain what makes a rule recursive.
  \item The arithmetic sequence $7,\; 11,\; 15,\; 19,\; \ldots$ can be
  written recursively as $u_{1} = 7$, $u_{n+1} = u_{n} + 4$, or explicitly
  as $u_{n} = 4n + 3$. Use each rule to find $u_{5}$. Then state which of
  the two rules you would use to find $u_{100}$, and why.
  \item The ninth to twelfth Fibonacci terms are $34,\; 55,\; 89,\; 144$.
  Use your calculator to find $\dfrac{u_{10}}{u_{9}}$ and
  $\dfrac{u_{12}}{u_{11}}$, and write down what you notice.
\end{enumerate}

\begin{tikzpicture}
  \draw (-0.6,0) rectangle (14.6,11.8);
  \foreach \y in {1,2,...,11} \draw[dotted] (0,\y)--(14.1,\y);
\end{tikzpicture}

\newpage

%--- Problem 5: Recursive self-improvement (RSI) in the AI news; arithmetic vs non-arithmetic growth [9 marks] ---
%--- Source: Original (freshly authored, AA SL 1.2); links to Problem 4 (recursion) ---
\item[5.]

In the news about artificial intelligence you will hear the phrase
\textbf{recursive self-improvement} (RSI): the idea that an AI system could
improve its own design, and that the improved system could then improve
itself again, round after round.

Here are two simple models of a system's capability score $c_{n}$ after
$n$ rounds. Both start at $c_{1} = 100$.

\begin{center}
\begin{tabular}{rl}
Model A: & each round adds $20$ points. \\[2pt]
Model B: & each round adds $20\%$ of the current score. \\
\end{tabular}
\end{center}

\begin{enumerate}[itemsep=0.1cm]
  \item Using your answer to Question 4(c), explain what the word
  \emph{recursive} is doing in the phrase ``recursive self-improvement''.
  \item Write down the first four terms of each model. State which model is
  arithmetic and give a reason.
  \item Write a recursive rule for each model.
  \item Find $c_{10}$ for each model.
  \item People who worry about RSI have in mind a process that speeds up as
  it goes. State which model matches that description, and explain what
  feature of the model makes it speed up.
\end{enumerate}

\begin{tikzpicture}
  \draw (-0.6,0) rectangle (14.6,10.5);
  \foreach \y in {1,2,...,9} \draw[dotted] (0,\y)--(14.1,\y);
\end{tikzpicture}

\newpage

%--- Problem 6: Arithmetic series calculations [9 marks] ---
%--- Source: Original (freshly authored, AA SL 1.2) ---
\item[6.]

Consider the arithmetic series
$$4 + 11 + 18 + \cdots + 200.$$

\begin{enumerate}[itemsep=0.1cm]
  \item Write down the value of the common difference $d$.
  \item Find the number of terms in the series.
  \item Find the sum of the series.
  \item An arithmetic sequence has $u_{1} = 6$ and $d = -2$. Find $S_{30}$.
  \item An arithmetic series has $u_{1} = 5$ and $S_{10} = 320$. Find the
  value of the common difference $d$.
\end{enumerate}

\begin{tikzpicture}
  \draw (-0.6,0) rectangle (14.6,14.5);
  \foreach \y in {1,2,...,13} \draw[dotted] (0,\y)--(14.1,\y);
\end{tikzpicture}

\end{enumerate}
\end{document}
